\documentclass[10pt,a4paper]{article}
\usepackage{graphicx}
\usepackage{array,booktabs,longtable}
\newcolumntype{P}[1]{>{\raggedright\arraybackslash}p{#1}}

\usepackage[margin=0.5in]{geometry}
\usepackage{amsmath,amssymb,bm}
\usepackage{siunitx}
\usepackage{booktabs,longtable,array}
\usepackage[hidelinks]{hyperref}
\usepackage{enumitem}
\usepackage{mathtools}

\title{Thermo-Richards-Mechanics (TRM) Model Overview}
\author{ChatGPT}
\date{}

\begin{document}

\maketitle

\section{Coupled Processes and Model Setup}

The provided OpenGeoSys input defines a Thermo-Richards-Mechanics (TRM) process -- a fully coupled thermal-hydraulic-mechanical model. In TRM, heat transport, variably-saturated groundwater flow (Richards' equation for a single liquid phase), and elastic deformation of the porous solid are solved \textbf{simultaneously} within one monolithic scheme \cite{ogs:trm, file:01_processes_TRM}. The \texttt{<process>} configuration in \texttt{01\_processes\_TRM.xml} confirms this: it specifies a single process of type \texttt{THERMO\_RICHARDS\_MECHANICS} named "TRM," with primary variables \textbf{displacement, pressure, and temperature} all in one process block \cite{file:01_processes_TRM}. This indicates a single fully-coupled process rather than separate sequential processes. A single Newton-Raphson solver (\texttt{basic\_newton}) is assigned to this process \cite{file:05_time_loop_TRM}, meaning all governing equations are solved together each time step (fully implicit coupling). There is no operator-splitting; the mechanical equilibrium, fluid flow, and heat transport equations are assembled into one system and iterated together (monolithic approach) \cite{ogs:trm}.

\subsection*{Processes included}
The TRM process couples: (1) \textbf{Heat transport} by conduction in the solid and liquid and convection by the liquid phase; (2) \textbf{Richards flow} (liquid water flow in unsaturated porous media) governed by Darcy's law; and (3) \textbf{Mechanical deformation} of the solid skeleton (linear elasticity). These couplings are configured via the input as follows:

\begin{itemize}
    \item \textbf{Energy (Thermal) Equation:} The energy balance for temperature $T$ includes an effective volumetric heat capacity and effective thermal conductivity of the medium \cite{ogs:trm}. Heat is conducted through both solid and fluid phases and can be carried by moving liquid. Mathematically, $(\rho c_p)^{\text{eff}} \dot{T} - \nabla\cdot(\mathbf{k}_T^{\text{eff}}\nabla T) + \rho^l c_p^l\,\mathbf{v}^l\cdot\nabla T = Q_T$ \cite{ogs:trm}, where $\rho c_p$ and $k_T$ denote density-heat capacity and thermal conductivity (with superscript eff for the combined medium). The input uses the \texttt{EffectiveThermalConductivityPorosityMixing} model for $\mathbf{k}_T^{\text{eff}}$ \cite{file:04_media_TRM} which automatically mixes the solid and liquid conductivities based on porosity and saturation \cite{ogs:trm}. The effective heat capacity $(\rho c_p)^{\text{eff}}$ is similarly computed from solid and liquid properties and the liquid saturation $S^l$ \cite{ogs:trm}. This means the solid's specific heat ($c_{p\_s}$) and density ($\rho_s$), and the liquid's specific heat ($c_{p\_l}$) and density, all contribute to the heat storage term in proportion to phase volumes \cite{ogs:trm}.

    \item \textbf{Richards (Hydraulic) Equation:} The liquid mass balance accounts for \textbf{storage/compressibility} and \textbf{unsaturated flow}. In the TRM formulation (single-phase Richards assumption), pore pressure $p^l$ is related inversely to capillary pressure ($p_c = -p^l$) \cite{ogs:trm}. The governing continuity equation (in a simplified form) is $S^l \beta\,\rho^l \dot{p^l} - \phi \frac{\partial S^l}{\partial p_c}\,\rho^l \dot{p^l} - S^l \rho^l(\alpha_B - S^l)\alpha_T^s \dot{T} + \nabla\cdot(\rho^l \mathbf{v}^l) + S^l \alpha_B\,\rho^l \nabla\cdot \dot{\mathbf{u}} = Q_H$ \cite{ogs:trm}. Here $\phi$ is porosity, $\beta$ represents effective fluid+solid compressibility, $\alpha_B$ is the Biot coefficient, and $\alpha_T^s$ is the solid's thermal expansivity. This equation shows several couplings: \textbf{fluid pressure} changes induce volumetric strain in the solid (via the Biot term $S^l\alpha_B\,\nabla\cdot\dot{\mathbf{u}}$) and vice versa; \textbf{temperature} changes affect fluid density and solid volumetric expansion ($\alpha_T^s$ term) which feed back into pore pressure \cite{ogs:trm}. The Darcy velocity for the liquid is given by $\mathbf{v}^l = -\frac{\mathbf{k}\,k_{\text{rel}}(S^l)}{\mu}\left(\nabla p^l - \rho^l \mathbf{g}\right)$ \cite{ogs:trm}. Thus, the \textbf{intrinsic permeability} tensor $\mathbf{k}$ (parameter $k_i$), the \textbf{relative permeability} $k_{\text{rel}}(S^l)$ (Van Genuchten model parameters), and the \textbf{dynamic viscosity} $\mu$ (parameter eta) govern fluid flow. In the input, $k_i$ is an anisotropic permeability (two principal values) assigned as a constant property \cite{file:04_media_TRM}, and a van Genuchten model is specified for relative permeability and saturation (with residual saturations and exponent) \cite{file:04_2_media_twophase}. Gravity $\mathbf{g}$ is turned off in this model (the \texttt{<specific\_body\_force>} is \texttt{0 0} in 2D, meaning no gravitational body force) \cite{file:01_processes_TRM}. As a result, the simulation does not impose a hydrostatic pressure gradient; the user instead provides an initial pressure distribution manually (see \texttt{ic\_pressure}) and boundary pressures.

    \item \textbf{Momentum (Mechanical) Equation:} The solid mechanics is formulated as a quasi-static force balance with poroelastic coupling: $\nabla\cdot[\boldsymbol{\sigma} - b(S^l)\,\alpha_B\,p^l \mathbf{I}] + \mathbf{f} = 0$ \cite{ogs:trm, doc:trm_gov_eq}. Here $\boldsymbol{\sigma}$ is the effective stress in the solid, $p^l$ the pore pressure, $\alpha_B$ (Biot coefficient) scales how much $p^l$ contributes to total stress, and $b(S^l)$ is a saturation-dependent Bishop's coefficient (effective stress partitioning in unsaturated conditions). The input uses a \textbf{Bishop's power-law model} with exponent 1 for $b(S)$ \cite{file:04_2_media_twophase}, which means $b(S^l)=S^l$ (pore pressure affects the stress in proportion to saturation). With Biot coefficient \texttt{biot = 1} \cite{file:03_parameters_TRM} at full saturation the pore pressure fully counteracts effective stress (Terzaghi's principle), while at lower $S^l$ the pore pressure influence is reduced ($b<1$). The primary unknown for mechanics is the \textbf{displacement} vector $\mathbf{u}$, which is related to strain $\boldsymbol{\epsilon}$ and stress via the linear elastic constitutive law: $\dot{\boldsymbol{\sigma}} = C : (\dot{\boldsymbol{\epsilon}} - \dot{\boldsymbol{\epsilon}}^T - \dot{\boldsymbol{\epsilon}}^{sw} \cdots)$ \cite{ogs:trm}. $C$ is the fourth-order stiffness tensor determined by the elastic moduli (Young's modulus, Poisson's ratio, etc.), $\dot{\boldsymbol{\epsilon}}^T$ is \textbf{thermal strain rate} (from thermal expansion of the solid), and $\dot{\boldsymbol{\epsilon}}^{sw}$ is \textbf{swelling strain} (strain from changes in saturation). In the input, the solid is defined as a linear orthotropic elastic material via two \texttt{<constitutive\_relation>} blocks (for two material regions) referencing parameters $E$, $G$, and $nu$ \cite{file:01_processes_TRM}. These give the Young's moduli, shear moduli, and Poisson ratios for the principal material axes (three values each, indicating orthotropic properties) \cite{file:03_parameters_TRM}. Thermal expansion of the solid is included by specifying a coefficient CTE (coefficient of thermal expansivity) as a solid-phase property \cite{file:04_media_TRM}. Additionally, a \textbf{saturation-dependent swelling} model is active: the \texttt{<property name="swelling\_stress\_rate">} in the solid phase uses \texttt{SaturationDependentSwelling} with specified parameters (swelling pressures and exponents) \cite{file:04_media_TRM, file:trm_medium_phase_prop}. This means that as the clay/rock becomes wet (increasing $S^l$), it undergoes an eigenstrain or stress increase (swelling) that is added to the mechanical equilibrium (modeled as an internal stress source term) -- effectively coupling moisture to mechanical strain beyond just pore pressure effects.
\end{itemize}

In summary, the TRM input files configure a \textbf{fully-coupled THM model} where heat, fluid, and mechanical processes influence each other. The coupling is \textbf{monolithic} (all primary variables -- temperature, pressure, displacement -- are solved together in one Newton iteration loop) \cite{file:05_time_loop_TRM}. For example, an increase in temperature can cause thermal expansion of the solid (via CTE) and thermal dilation of pore water (via fluid density's temperature dependence), affecting stress and pore pressure \cite{ogs:trm}. Likewise, a drop in pore pressure (e.g., due to drainage at a boundary) leads to effective stress increase and deformation (governed by Biot's coefficient and Bishop's $b(S)$) \cite{ogs:trm, doc:trm_gov_eq}, and conversely compression of the solid can expel water and change pressure (the $\nabla\cdot \dot{\mathbf{u}}$ term in the fluid mass balance) \cite{ogs:trm}. All these interactions are captured by the single TRM process as configured.

The table below and following sections detail all \textbf{input parameters} used in the model, including their types, locations in the XML, and their roles in the governing equations.

\section{Material and Input Parameters}

\subsection*{Overview}
The model's input parameters can be grouped into (a) initial conditions, (b) boundary conditions, and (c) material properties (covering porous medium, solid phase, and liquid phase properties). The parameters are defined mostly in \texttt{03\_parameters\_TRM.xml} (using the \texttt{<parameter>} tag) and referenced elsewhere, while some are given directly in the \texttt{04\_media\_TRM.xml} property definitions. Table 1 and Table 2 below list all parameters, along with units (where available or physically relevant), where they appear in the input, and what they represent in the model's physics.

\subsection*{Initial Conditions}
\begin{itemize}
    \item \textbf{ic\_displacement} -- Initial displacement, set as a Constant vector \texttt{[0 0]} in \texttt{03\_parameters\_TRM.xml} \cite{file:03_parameters_TRM}. Units: meters (m). This provides the initial condition for the solid displacement field (here zero, meaning no initial deformation). It is referenced in the \texttt{<initial\_condition>} of the displacement process variable \cite{file:02_process_variables_TRM}. Role: serves as $\mathbf{u}(t=0)$, the starting displacement, typically zero in consolidation problems.

    \item \textbf{ic\_pressure} -- Initial pore pressure distribution, defined as a Function in \texttt{03\_parameters\_TRM.xml} \cite{file:03_parameters_TRM}. Units: Pascals (Pa). The given expression is \texttt{1.5e6 - ((1095*(1 + 3.2e-10*(1.5e6 - 7.5e6))))*0*y}, which essentially evaluates to $1.5 \times 10^6$ Pa uniformly (because it multiplies by \texttt{0*y}, nullifying the intended gradient). This was likely meant to represent an initial hydrostatic pressure profile with a small gradient (using water density $\sim$1095 kg/m³ and compressibility 3.2e-10 Pa\textsuperscript{-1}), but with the coefficient set to 0, it simplifies to a constant 1.5 MPa. It is applied as the initial condition for the pressure field \cite{file:02_process_variables_TRM}. Role: $p^l(t=0)$ in the Richards equation. Here it initializes the pore pressure to $\sim$1.5 MPa everywhere (which, given gravity is off, is an equilibrium state). The mention of 1095 and 3.2e-10 in the formula aligns with the liquid's density model (reference density and compressibility), indicating an attempt to be consistent with fluid properties.

    \item \textbf{ic\_temperature} -- Initial temperature, Constant 298.15 K (25 °C) \cite{file:03_parameters_TRM}. This sets the initial temperature of the domain uniformly to 298.15 K via the temperature process variable's initial condition \cite{file:02_process_variables_TRM}. It serves as $T(t=0)$ for the heat transport equation. Units: Kelvin (K). With this and a fixed boundary temperature (see below), the thermal field starts and stays at 25 °C unless other influences (not present here) cause changes.

    \item \textbf{ic\_sigma0} -- Initial stress state, given as a Function with multiple components \cite{file:03_parameters_TRM}. Units: Pascals. This parameter provides the initial stress tensor components in the solid (likely $\sigma_{xx}, \sigma_{yy}, \sigma_{zz}, \tau_{xy}, \ldots$). The expressions suggest an initial vertical stress of about -7.0e6 Pa and horizontal stress of -1.21e6 Pa (negative sign indicating compressive stress) with some minor terms (the \texttt{0*y} again nullifies any gradient). However, \textbf{this initial stress parameter is not actually applied in the model} because the \texttt{<initial\_stress>} tag in the process was commented out \cite{file:01_processes_TRM}. So, `ic\_sigma0` is defined but not used (no initial stress is set in the mechanical process). If it were used, it would supply the in-situ stress field for the solid, providing $\boldsymbol{\sigma}(t=0)$ in the momentum equation (important in geomechanical problems to account for pre-existing stress). In this case, since it's omitted, the initial stress defaults to zero (and the boundaries are constrained to prevent sudden imbalance).
\end{itemize}

\subsection*{Boundary Conditions}
\begin{itemize}
    \item \textbf{bc\_displacement} -- Displacement fixed-value boundary, Constant 0 (m) \cite{file:03_parameters_TRM}. It is applied as a Dirichlet boundary condition on multiple boundary meshes (\texttt{cd-a\_left}, \texttt{cd-a\_right}, \texttt{cd-a\_top}, \texttt{cd-a\_bottom}) for specified displacement components \cite{file:02_process_variables_TRM}. Essentially, this means the left, right, top, and bottom edges of the model are fixed (zero displacement) in the normal directions. Physically, this simulates supported or symmetry boundaries with no movement. Role in PDE: it sets $\mathbf{u}(t) = 0$ on those boundaries, providing mechanical constraints so that the solid can't deform there.

    \item \textbf{bc\_pressure\_outside} -- Fixed pore pressure boundary, Constant $1.5 \times 10^6$ Pa \cite{file:03_parameters_TRM}. (This likely represents an external reservoir pressure of 1.5 MPa.) It is intended for a boundary condition, but in the provided process variables, the only pressure Dirichlet BCs used are \texttt{bc\_top\_pressure} and \texttt{open\_niche\_seasonal} (see below) -- there is \textbf{no reference} to \texttt{bc\_pressure\_outside} in the \texttt{<boundary\_condition>} list for pressure \cite{file:02_process_variables_TRM, file:trm_proc_var_bc}. This suggests \texttt{bc\_pressure\_outside} was defined but not actually applied to any boundary (perhaps an artifact of a previous setup). If it were used, it would enforce 1.5 MPa on a boundary (likely intended for the outer boundary "outside" of the model domain). Since it's unused, it does not directly affect the model.

    \item \textbf{bc\_top\_pressure} -- Top boundary pressure, given as a Function \cite{file:03_parameters_TRM}. Units: Pa. The expression is \texttt{1.5e6 - ((1095*(1 + 3.2e-10*(1.5e6 - 7.5e6))))*0*50}, which again simplifies to 1.5e6 Pa (because of \texttt{0*50}). This was likely meant to calculate pressure at the top boundary (perhaps at 50 m above a reference elevation) considering fluid weight, but as written it yields 1.5 MPa flat. It is applied as a Dirichlet condition on the top boundary mesh (\texttt{cd-a\_top}) of the model \cite{file:02_process_variables_TRM}. Thus, the top boundary's pore pressure is fixed at $\sim$1.5 MPa (the same as the initial pressure). This essentially maintains the initial pressure at the top, preventing outflow or inflow through the top boundary initially. Physical role: simulates a constant pressure boundary (perhaps connected to a water reservoir or an assumed constant head at the top of the formation).

    \item \textbf{open\_niche\_seasonal} -- Niche pressure boundary (time-variable), defined as a \texttt{CurveScaled} parameter in \texttt{03\_parameters\_TRM.xml} \cite{file:03_parameters_TRM}. Units: Pa (assumed). This uses a time \textbf{curve} (\texttt{open\_niche\_seasonal\_curve}) scaled by a factor (\texttt{pressure\_scaling\_factor} = 1) to produce a time-dependent pressure value. It is applied as a Dirichlet boundary condition on the mesh \texttt{cd-a\_niche4} \cite{file:02_process_variables_TRM}, which presumably represents an open excavation (a niche or gallery opening in the rock). In effect, this boundary is exposed to atmospheric or seasonal pressure variations. The actual numeric values come from the curve data (not visible in the snippet, but likely the curve represents seasonal cycles of air pressure or humidity influence). Given the context, it likely oscillates around atmospheric pressure. Additionally, \texttt{bc\_atm\_pressure} (0.1 MPa) is defined as a constant reference for atmospheric pressure \cite{file:03_parameters_TRM}, though it's not directly used by the boundary condition. Role: This provides a time-varying pore pressure on the niche boundary, simulating, for example, seasonal ventilation or atmospheric pressure changes in an open tunnel. In the Richards equation, this is a Dirichlet condition $p^l(t)$ on that boundary, driving flow if the interior pressure differs from this value.

    \item \textbf{bc\_temperature\_outside} -- Fixed temperature boundary, Constant 298.15 K (25 °C) \cite{file:03_parameters_TRM}. It is applied to the mesh \texttt{bulk\_all} as a Dirichlet condition for temperature \cite{file:02_process_variables_TRM}. If \texttt{bulk\_all} represents the entire external boundary of the domain, this essentially fixes all outer boundaries at 25 °C. (It might even have been intended as an initial condition if misused, but the tag is under \texttt{boundary\_conditions}.) With the initial temperature also at 25 °C, this means the domain's boundaries are held at the same temperature as the interior initially, resulting in no thermal gradients. Role: ensures a constant temperature at the boundary (an isothermal boundary at 25 °C), so any internal thermal perturbation would dissipate to this temperature. In practice, since no heat source is present and initial T equals boundary T, the thermal field remains uniform.

    \item \textbf{load\_top} -- Mechanical surface load, defined as a Function (expression \texttt{(-7e6) + ... *0*50}) in the parameters file \cite{file:03_parameters_TRM}. Units: Pa (as a stress or pressure applied). This was intended as a \textbf{Neumann boundary condition} (surface traction) on the top, representing perhaps an overburden stress of 7 MPa (with some minor adjustment, again nullified by \texttt{0*50}). However, in \texttt{02\_process\_variables\_TRM.xml}, the corresponding boundary condition is commented out (the snippet shows a commented \texttt{<boundary\_condition>} for \texttt{load\_top} on \texttt{cd-a\_top} as a Neumann load) \cite{file:02_process_variables_TRM}. Thus, \texttt{load\_top} is \textbf{not actually applied} in the simulation. If it were active, it would impose a downward normal stress of $\sim$7 MPa on the top boundary (simulating weight of overlying rock or a pressure chamber). Since it's not used, the top is instead fixed (Dirichlet) in displacement, and no external force is applied. This parameter can be disregarded for the running model, though it indicates the user considered applying an external load equal to the in-situ stress.

    \item \textbf{bc\_atm\_pressure} -- Atmospheric pressure reference, Constant $0.1 \times 10^6$ Pa (0.1 MPa) \cite{file:03_parameters_TRM}, equivalent to $\sim$1 bar. This is a parameter for atmospheric pressure. It is not directly referenced in the active boundary conditions (the open niche uses its own curve, not this constant). It might have been intended for use in a "closed niche" scenario or for comparison. In this model configuration, it serves as a documented value of atmospheric pressure but does not enter the equations unless the curve or other logic uses it (which in the given files, they do not). We mention it for completeness, but it has no effect unless manually used.

    \item \textbf{pressure\_scaling\_factor} -- Constant 1.0 (dimensionless) \cite{file:03_parameters_TRM}. This parameter is used as the base parameter in the \texttt{CurveScaled} definition of \texttt{open\_niche\_seasonal} \cite{file:03_parameters_TRM}. Essentially, \texttt{open\_niche\_seasonal} = (\texttt{pressure\_scaling\_factor}) $\times$ (value from \texttt{open\_niche\_seasonal\_curve}). Since it's 1.0, the curve values are taken as-is. This indirection allows easily changing the amplitude of the curve by adjusting one number. Physical role: not a physical quantity itself, but part of how the time-dependent boundary condition is constructed. (If the curve was normalized, one could set this to 0.1e6 to scale it to Pascals, but here it's 1.0, implying the curve likely provides absolute pressure in Pa.)
\end{itemize}

% Table 1
% Table 1 (spans multiple pages)
\setlength{\tabcolsep}{4pt}         % tighter columns (optional)
\renewcommand{\arraystretch}{1.12}  % a bit more row height (optional)

\begin{longtable}{P{4cm} P{1.9cm} P{0.9cm} P{4cm} P{6cm}}
\caption{Initial and Boundary Condition Parameters (OpenGeoSys inputs)}\label{tab:ic-bc-params}\\
\toprule
\textbf{Name} & \textbf{Type} & \textbf{Units} & \textbf{Input Location} & \textbf{Physical Role in Model} \\
\midrule
\endfirsthead

\multicolumn{5}{l}{\small\itshape Table~\thetable{} (continued)}\\
\toprule
\textbf{Name} & \textbf{Type} & \textbf{Units} & \textbf{Input Location} & \textbf{Physical Role in Model} \\
\midrule
\endhead

\midrule
\multicolumn{5}{r}{\small\itshape Continued on next page}\\
\endfoot

\bottomrule
\endlastfoot

\texttt{ic\_displacement} & Constant & m & \texttt{03\_parameters\_TRM.xml}, \texttt{<parameter>} \cite{file:03_parameters_TRM} & Initial displacement (vector) for solid; here [0,0] (no initial deformation). Used in \texttt{<initial\_condition>} for displacement DOF \cite{file:02_process_variables_TRM}. \\
\texttt{ic\_pressure} & Function & Pa & \texttt{03\_parameters\_TRM.xml}, \texttt{<parameter>} \cite{file:03_parameters_TRM} & Initial pore pressure field. Expression given (intended hydrostatic profile) simplifies to 1.5e6 Pa uniform. Sets initial pressure for Richards flow \cite{file:02_process_variables_TRM}. \\
\texttt{ic\_temperature} & Constant & K & \texttt{03\_parameters\_TRM.xml}, \texttt{<parameter>} \cite{file:03_parameters_TRM} & Initial temperature of domain. Here 298.15 K everywhere (25 °C). Sets initial T for heat transport \cite{file:02_process_variables_TRM}. \\
\texttt{ic\_sigma0} & Function & Pa & \texttt{03\_parameters\_TRM.xml}, \texttt{<parameter>} \cite{file:03_parameters_TRM} & Initial stress tensor in solid (expressions for components). \textbf{Defined but not applied} (\texttt{initial\_stress} tag commented out \cite{file:01_processes_TRM}); would provide in-situ stress if used. \\
\texttt{bc\_displacement} & Constant & m & \texttt{03\_parameters\_TRM.xml}, \texttt{<parameter>} \cite{file:03_parameters_TRM} & Displacement Dirichlet BC = 0 (fixed support). Applied to model boundaries (left/right/top/bottom) for specified components \cite{file:02_process_variables_TRM}, fixing those boundaries in place (no movement). \\
\texttt{bc\_pressure\_outside} & Constant & Pa & \texttt{03\_parameters\_TRM.xml}, \texttt{<parameter>} \cite{file:03_parameters_TRM} & Fixed pressure = 1.5e6 Pa. \textbf{Defined but not used} in this model's BC list (likely intended for an outer boundary). No direct effect since not referenced \cite{file:trm_proc_var_bc}. \\
\texttt{bc\_top\_pressure} & Function & Pa & \texttt{03\_parameters\_TRM.xml}, \texttt{<parameter>} \cite{file:03_parameters_TRM} & Top boundary pressure Dirichlet. Expression evaluates to $\sim$1.5e6 Pa constant (no gradient). Applied on \texttt{cd-a\_top} mesh \cite{file:02_process_variables_TRM}. Represents a constant pressure (e.g., connected to a reservoir or overburden) at the top boundary. \\
\texttt{bc\_temperature\_outside} & Constant & K & \texttt{03\_parameters\_TRM.xml}, \texttt{<parameter>} \cite{file:03_parameters_TRM} & Fixed temperature = 298.15 K on external boundary (\texttt{bulk\_all}) \cite{file:trm_proc_var_bc}. Maintains 25 °C at boundaries (isothermal boundary condition). \\
\texttt{load\_top} & Function & Pa & \texttt{03\_parameters\_TRM.xml}, \texttt{<parameter>} \cite{file:03_parameters_TRM} & Top surface mechanical load (Neumann BC) expression ($\sim$-7e6 Pa). \textbf{Not active} (the Neumann BC is commented out) so no force applied \cite{file:02_process_variables_TRM}. Included for completeness – would simulate a 7 MPa downward stress if enabled. \\
\texttt{bc\_atm\_pressure} & Constant & Pa & \texttt{03\_parameters\_TRM.xml}, \texttt{<parameter>} \cite{file:03_parameters_TRM} & Atmospheric pressure = 0.1e6 Pa (1 bar). Not directly used in BCs (the open niche uses a curve instead). Serves as a reference value (could be used for a closed boundary scenario or in the curve). No effect on model as is. \\
\texttt{pressure\_scaling\_factor} & Constant & - & \texttt{03\_parameters\_TRM.xml}, \texttt{<parameter>} \cite{file:03_parameters_TRM} & Scaling factor = 1.0 for the \texttt{open\_niche\_seasonal} curve. Unitless. Multiplies the curve's values (identity scaling here). Used in CurveScaled BC definition \cite{file:03_parameters_TRM} to allow easy re-scaling of niche pressure if needed. \\
\texttt{open\_niche\_seasonal} & CurveScaled & Pa (output) & \texttt{03\_parameters\_TRM.xml}, \texttt{<parameter>} \cite{file:03_parameters_TRM} & Time-dependent pressure for open niche boundary. Defined by a time curve (\texttt{open\_niche\_seasonal\_curve}) scaled by factor (1.0). Applied as Dirichlet BC on \texttt{cd-a\_niche4} (the niche interior) \cite{file:02_process_variables_TRM}. Represents seasonal variation of air/fluid pressure in the open gallery. \\
\end{longtable}


\subsection*{Porous Medium \& Solid Material Properties}

% Table 2 - Part 1 (as running text for better flow)
\begin{itemize}
    \item \textbf{phi (porosity)} - Constant 0.105 (dimensionless) \cite{file:03_parameters_TRM}. Defined in \texttt{03\_parameters\_TRM.xml} and used as the \texttt{<property name="porosity" type="Parameter">} in the medium definition \cite{file:04_media_TRM}. Porosity is the fraction of volume that is pore space. Role in model: appears in the Richards equation and thermal calculations. It multiplies saturation in storage terms and in the effective heat capacity: e.g. $(1-\phi)\rho^s c_p^s + \phi S^l \rho^l c_p^l$ \cite{ogs:trm}. A smaller porosity (0.105) means relatively little pore volume (typical of rock), affecting how much fluid can be stored and how much the fluid's heat capacity contributes.

    \item \textbf{k\_i (intrinsic permeability)} – Constant tensor, values [1.00E-19, 0, 0, 0.40E-19] (m²) in 2D \cite{file:03_parameters_TRM}. This represents an \textbf{anisotropic permeability}: $k_{xx}=1.0\times10^{-19}$ m² and $k_{yy}=0.4\times10^{-19}$ m², with off-diagonals 0 (no cross-anisotropy). It's applied via the \texttt{<property name="permeability" type="Parameter">} in the medium \cite{file:04_media_TRM}. Role: $\mathbf{k}$ in Darcy's law \cite{ogs:trm}, controlling how easily fluid flows through the porous matrix. Such a low permeability ($\sim 10^{-19}$) is characteristic of very tight rock/clay. The anisotropy (factor $\sim$0.4 in vertical direction) indicates the vertical permeability is 40% of horizontal, which will influence the flow pattern (flow is easier along the higher permeability direction). Units are m² (intrinsic permeability).

    \item \textbf{biot (Biot coefficient)} – Constant 1 (dimensionless) \cite{file:03_parameters_TRM}. Set as \texttt{<property name="biot\_coefficient" type="Parameter">} in the medium \cite{file:04_media_TRM}. Biot's coefficient $\alpha_B$ represents how much pore pressure contributes to volumetric stress in the solid (1 means fluid pressure fully couples to stress, typical for soft porous materials or full saturation assumption). Role: appears in the effective stress law and fluid mass balance. In the momentum equation, it multiplies $p^l$ (along with the Bishop factor) in the term $\alpha_B b(S) p^l$ \cite{ogs:trm, doc:trm_gov_eq}, effectively subtracting that from total stress to get effective stress. In the fluid equation, $\alpha_B$ multiplies the volumetric strain term $\nabla\cdot \dot{\mathbf{u}}$ \cite{ogs:trm}, meaning any volume change of the skeleton contributes to pore volume change. With $\alpha_B=1$, the model assumes all solid volume change translates to pore volume change (no matrix compression decoupling).

    \item \textbf{E (Young's modulus)} – Constant (vector) = [13.80e+9, 6.00e+9, 13.80e+9] Pa \cite{file:03_parameters_TRM}. This gives the Young's moduli in the principal material directions (X, Y, Z). The values indicate the rock is stiffer in X and Z directions (13.8 GPa) and softer in Y (6.0 GPa), suggesting an orthotropic elasticity (perhaps bedding or layering makes the vertical direction different). These are used in the \texttt{<constitutive\_relation LinearElasticOrthotropic>} for the solid's mechanical law \cite{file:01_processes_TRM}. Role: Along with G and nu, determines the elastic stress-strain relationship $C_{ijkl}$. Specifically, $E$ controls axial stiffness in each principal direction. A high $E$ (13.8 GPa) means the material is quite stiff in those directions, whereas 6.0 GPa in the weaker direction means it's more compressible/flexible in that orientation. In the generalized Hooke's law, these define how much stress results from a given strain (and vice versa). They affect deformation under load and coupling via the term $C:\dot{\boldsymbol{\epsilon}}^e$ in the stress equation \cite{ogs:trm}.

    \item \textbf{G (Shear modulus)} – Constant (vector) = [4.79e+9, 2.46e+9, 4.79e+9] Pa \cite{file:03_parameters_TRM}. These are the shear moduli corresponding to the material's principal directions (for shear in planes orthogonal to each principal axis). The values again show anisotropy (shear modulus is lowest for shear involving the Y-direction). In an orthotropic linear elastic model, $G$ (along with $E$ and $\nu$) defines off-diagonal stiffness terms. Role: Shear modulus governs the material's resistance to shear deformation. In the elasticity tensor $C$, $G_{ij}$ appears in the relations for shear stress-strain. Higher $G$ (4.79 GPa) in XZ plane means more resistance to shear in those planes than in the plane involving Y (2.46 GPa). These moduli will affect how the rock responds to shear strains induced by e.g. uneven stresses or swelling.

    \item \textbf{nu (Poisson's ratios)} – Constant (vector) = [0.44, 0.22, 0.44] (dimensionless) \cite{file:03_parameters_TRM}. These are the Poisson's ratios $\nu_{xy}, \nu_{yz}, \nu_{zx}$ (or analogous, depending on convention) for the orthotropic material. They indicate the coupling between axial and lateral strain in different directions. Here, 0.44 in two directions and 0.22 in one suggests that in two planes the material is highly Poisson-active (0.44 is quite high, meaning if you stretch in one direction it contracts 44\% as much in the perpendicular direction), whereas in the other it's about 0.22. Role: Enter the elastic constitutive law -- they influence the $C$ matrix components that couple normal strains. For example, a high $\nu$ (0.44) might represent near-incompressibility in one plane (or a material that is constrained such that it cannot expand laterally much when compressed vertically). Combined with $E$, these determine the bulk compressibility of the solid matrix (which in turn affects the “$\beta$” composite compressibility in the fluid equation \cite{ogs:trm}, although here Biot=1 and these $\nu$ values mostly affect how stress translates to strain rather than directly entering fluid storage, since fluid storage is handled by the Beta term and porosity).

    \item \textbf{rho\_s (solid grain density)} – Constant 2520 kg/m³ \cite{file:03_parameters_TRM}. Used in the solid phase \texttt{<property name="density" type="Parameter">} \cite{file:04_media_TRM}. This is the density of the solid matrix (mineral density of the rock). Role: contributes to the \textbf{effective heat capacity} of the porous medium and can contribute to gravitational body force if used. In the energy equation, solid density appears in $(\rho c_p)^{\text{eff}}$ \cite{ogs:trm}. Also, if body forces were enabled, $\rho_s$ with gravity would give solid weight per volume (but here specific\_body\_force is zero, so no direct use in momentum). The chosen value 2520 kg/m³ suggests a typical rock (e.g., sandstone or claystone).

    \item \textbf{kappa\_th\_s (solid thermal conductivity)} – Constant 2.156 W/(m·K) \cite{file:03_parameters_TRM}. Assigned to solid phase \texttt{<property name="thermal\_conductivity">} \cite{file:04_media_TRM}. It represents how well the solid phase conducts heat. Role: in the \texttt{EffectiveThermalConductivityPorosityMixing} model, the solid's thermal conductivity $\mathbf{k}_T^s$ is combined with the liquid's ($\kappa_l$) according to $k_T^{\text{eff}} = (1-\phi) k_T^s + \phi S^l k_T^l$ \cite{ogs:trm}. Thus, this value (2.156 W/m·K) will dominate heat conduction when pores are mostly solid (low porosity or dry conditions). It's in a typical range for rocks (e.g., claystone might be $\sim$2 W/mK).

    \item \textbf{c\_p\_s (solid specific heat capacity)} – Constant 1254.74 J/(kg·K) \cite{file:03_parameters_TRM}. Assigned to solid phase \texttt{<property name="specific\_heat\_capacity">} \cite{file:04_media_TRM}. This is the heat capacity of the solid grains. Role: appears in the effective volumetric heat capacity $(\rho c_p)^{\text{eff}} = (1-\phi)\rho^s c_p^s + \phi S^l \rho^l c_p^l$ \cite{ogs:trm}. So 1254.74 J/kgK combined with 2520 kg/m³ gives the solid's volumetric heat capacity contribution. It affects how much heat the solid phase can store per degree temperature change. This value is within expected range for rock minerals.

    \item \textbf{CTE (coefficient of thermal expansivity)} – Constant 1254.74 (units in input comment are mistakenly [J/kgK], but physically this should be 1/K) \cite{file:03_parameters_TRM}. This parameter is used in solid phase \texttt{<property name="thermal\_expansivity">} \cite{file:04_media_TRM}. It represents the linear thermal expansion coefficient $\alpha_T^s$ of the solid skeleton. The numerical value 1254.74 is \textbf{unusually high} for a thermal expansion (likely a unit error or placeholder - typical $\alpha_T^s$ for rock is on the order of 1e-5 to 1e-6 1/K). It might be that the user intended $1.254\times10^{-5}$ K\textsuperscript{-1} and mis-entered a scaled value. In any case, in the model this $\alpha_T^s$ enters the thermal strain term in the stress equation and the thermal pressurization term in the fluid equation. Specifically, thermal strain $\boldsymbol{\epsilon}^T = \alpha_T^s \Delta T \mathbf{I}$ contributes to reducing effective stress when temperature rises (if the solid can expand) \cite{ogs:trm}. In the mass balance, $\rho^l(\alpha_B - S^l)\alpha_T^s \dot{T}$ is a term \cite{ogs:trm}, meaning that if the solid heats up (and tries to expand) while pores are not fully saturated ($S^l < \alpha_B$), it will expel some fluid or create pressure. Thus, CTE couples temperature to both mechanical and hydraulic responses. (Given the odd value, this may overestimate thermal strain unless the value was meant in another unit.)

    \item \textbf{Swelling model parameters:} The solid has a \texttt{swelling stress rate} property of type \texttt{SaturationDependentSwelling} \cite{file:04_media_TRM, file:trm_medium_phase_prop}, which introduces additional stress as saturation rises (simulating e.g. clay swelling). Key parameters inside this model:
    \begin{itemize}
        \item \texttt{swelling\_pressures} - Three values [1.0E+06, 1.0E+06, 1.0E+06] Pa \cite{file:04_media_TRM}. Likely the swelling pressure coefficients in the three principal directions (here equal, implying isotropic swelling). On full saturation, the model will generate a stress of up to $\sim$1 MPa in each principal direction (if we assume linear relation) resisting the swelling.
        \item \texttt{exponents} - [1, 1, 1] \cite{file:04_media_TRM}. These exponents define the nonlinear relation between saturation and swelling stress (all 1 means linear proportionality: swelling stress $\propto S^l$ between the limits).
        \item \texttt{lower\_saturation\_limit} - 0.1 (dimensionless saturation) \cite{file:04_media_TRM}. Below 0.1 saturation, swelling is not generated (dry state – no swelling).
        \item \texttt{upper\_saturation\_limit} - 1.0 \cite{file:04_media_TRM}. At full saturation (100\% water), the model applies the full swelling pressure. Between 10\% and 100\% saturation, the swelling stress increases (linearly in this case) from 0 to 1 MPa.
        \item \textbf{Role:} These parameters together mean that once the pores start filling beyond 10\% water, the solid experiences an internal stress pushing outward, up to 1 MPa at full saturation. This effectively adds a source term in the stress equilibrium (swelling strain $\dot{\boldsymbol{\epsilon}}^{sw}$ in the constitutive law \cite{ogs:trm}) that is coupled to moisture. It does not correspond to a classic material property like elasticity or thermal expansion, but rather a user-defined empirical way to model phenomena like clay swelling pressure. In the mechanical equation, it will appear as an added stress (like an expansion strain) that must be counteracted by constraints or cause deformation if free. In summary, it couples the increase of $S^l$ to an increase in stress, which in a constrained system could lead to pressure increase or in an unconstrained system to expansion.
    \end{itemize}
\end{itemize}

\subsection*{Liquid (Fluid) Properties}
\begin{itemize}
    \item \textbf{c\_p\_l (liquid specific heat capacity)} – Constant 4160 J/(kg·K) \cite{file:03_parameters_TRM}. This is the specific heat of water (or brine) and is used in the liquid phase \texttt{<property name="specific\_heat\_capacity">} \cite{file:04_1_media_aqu_liq}. Role: contributes to the energy balance. In the effective heat capacity, $\phi S^l \rho^l c_p^l$ term uses this value \cite{ogs:trm}, and in the advective heat term $\rho^l c_p^l \mathbf{v}^l \cdot \nabla T$ \cite{ogs:trm}, it multiplies how much heat is carried by flowing water. 4160 J/kgK is slightly higher than pure water (which is $\sim$4182 J/kgK at 20°C), but in the correct ballpark.

    \item \textbf{kappa\_l (liquid thermal conductivity)} – Constant 0.6 W/(m·K) \cite{file:03_parameters_TRM}. Assigned to liquid phase \texttt{<property name="thermal\_conductivity">} \cite{file:04_1_media_aqu_liq}. Represents thermal conductivity of water. Role: used in the effective thermal conductivity mixing: $k_T^{\text{eff}} = (1-\phi)k_T^s + \phi S^l k_T^l$ \cite{ogs:trm}. Water's thermal conductivity ($\sim$0.6 W/mK) is lower than the solid's (2.156), so as saturation increases, the bulk thermal conductivity will shift towards a weighted average. This affects heat diffusion through wet vs dry porous media.

    \item \textbf{eta (dynamic viscosity)} – Constant 1.0e-3 Pa·s \cite{file:03_parameters_TRM}. Used in liquid phase \texttt{<property name="viscosity">} \cite{file:04_1_media_aqu_liq}. This is the viscosity of water ($\sim$1 cP) at room temperature. Role: appears in Darcy's law denominator $\mathbf{v}^l = -\frac{\mathbf{k} k_{\text{rel}}}{\mu}(\nabla p - \rho^l \mathbf{g})$ \cite{ogs:trm}. A higher viscosity would resist flow, but 1e-3 Pa·s is standard for water at 20°C. Since the model's temperature range is near 25°C, assuming constant viscosity is reasonable. This parameter thus controls the ease of fluid flow for a given pressure gradient and permeability. (Note: the input comment mistakenly had "[J/kgK]" for units, which is incorrect; the correct unit is Pa·s.)

    \item \textbf{Liquid density model:} Instead of a fixed parameter, the liquid phase density is given by a functional \textbf{Linear model} in \texttt{04\_1\_media\_aqu\_liq.xml} \cite{file:04_1_media_aqu_liq}:
    \begin{itemize}
        \item \textbf{Reference density:} 1095 (kg/m³) at reference conditions (reference temperature 298.15 K and reference pressure 7.5e6 Pa) \cite{file:04_1_media_aqu_liq}. This suggests that at 25 °C and 7.5 MPa, the water has density 1095 kg/m³ - slightly higher than pure water (perhaps due to dissolved salts or simply a value chosen for deep groundwater at pressure).
        \item \textbf{Thermal expansion slope:} -5$\times 10^{-4}$ per K \cite{file:04_1_media_aqu_liq}. This means for each 1 K increase in temperature, the density decreases by 0.0005$\times$(reference density) or roughly 0.05\% per K. In absolute terms, $\partial \rho / \partial T = -0.5475$ kg/m³K (at ref conditions). This is a mild thermal expansivity of water – e.g., from 25°C to 26°C, density would drop slightly. Role: this slope contributes to the term $\partial \rho^l/\partial T$ in the mass balance \cite{ogs:trm}. A positive temperature change will reduce $\rho^l$, tending to lower fluid mass or pressure unless constrained. Essentially, heating causes slight de-pressurization if volume is constant (thermal pressurization is also influenced by the solid's expansion via $\alpha_T^s$ term).
        \item \textbf{Compressibility slope:} 3.2$\times 10^{-10}$ per Pa \cite{file:04_1_media_aqu_liq}. This is $\partial \rho / \partial p$ (1/Pa). Multiply by a pressure change to get density change. For example, for a 1 MPa increase, density increases by 3.2e-4 of reference, which is $\sim$0.35 kg/m³. This corresponds to a water compressibility of about 3.2e-10 Pa\textsuperscript{-1} (which is in the same order as typical water compressibility $\sim$4.4e-10 Pa\textsuperscript{-1}). Role: this defines how much density (and thus fluid mass) increases with pressure - contributing to the \textbf{storage term} of the fluid. In the governing equation, the composite $S^l \beta$ includes liquid compressibility; effectively $\beta \approx (1/\rho^l) \partial \rho^l/\partial p$ plus any solid matrix compressibility. Here, with $\partial \rho / \partial p$ given, the model can compute changes in fluid content with pressure. This affects how pressure waves travel (e.g., higher compressibility means more storativity, slower pressure changes). Also, in Darcy's law, the $\rho^l \mathbf{g}$ term uses density which will be slightly larger at high pressure (though gravity is off, so $\rho g$ term is moot here).
        \item These values are built into the model; thus, there is \textbf{no explicit} \texttt{rho\_l} parameter in the parameters file. The physical quantity is fluid density $\rho^l(T,p)$. At 25°C and 1.5 MPa (top boundary), evaluating the linear law: starting from 1095 at 7.5 MPa, reducing pressure by 6 MPa (to 1.5 MPa) would lower density by $6\text{e}6 \times 3.2\text{e}-10 = 1.92\text{e}-3$ fraction ($\sim$2 kg/m³), and the thermal term is zero if T=298.15 reference. So at 1.5 MPa, 25°C, $\rho^l \approx 1093$ kg/m³. At the top vs bottom of domain the difference is marginal (the input formula comment noted "just 1\% difference if not considering change in rho" \cite{file:03_parameters_TRM}). Role in model: $\rho^l$ enters the continuity equation (fluid mass accumulation $\propto \rho^l$) and Darcy flux (via gravity term and slight effect on $\nabla p$ if solving in terms of head). In our case with no gravity and nearly constant p, the compressibility mainly contributes to a very small storage effect.
    \end{itemize}

    \item \textbf{Unsaturated flow (Van Genuchten) parameters:} The medium properties include a van Genuchten type capillary pressure-saturation and relative permeability model \cite{file:04_2_media_twophase}:
    \begin{itemize}
        \item \textbf{residual\_liquid\_saturation ($S_{lr}$)} – 0.1 (-) \cite{file:04_2_media_twophase}. This means 10\% of the pore space remains filled with water even at "dry" conditions (and conversely, residual gas saturation is 0 in this model, implying if fully saturated, no trapped gas). Role: in the \textbf{Saturation-capillary pressure relation}, below $S^l = 0.1$ the model treats the liquid as effectively non-wetting residual. It appears in the Van Genuchten formulas for $S_e$ (effective saturation) and influences relative permeability. In OGS input, it's set for both relative\_permeability and saturation models.
        \item \textbf{residual\_gas\_saturation ($S_{gr}$)} – 0 (-) \cite{file:04_2_media_twophase}. No residual gas, meaning at $S^l=1$ the medium is fully saturated with no trapped air. This simplifies the Bishop's $b(S)$ at full saturation to 1. Role: used in Van Genuchten relationships; here essentially $S_e = (S^l - 0.1) / (1-0.1-0)$.
        \item \textbf{exponent (m or related to n)} – 0.45 (-) \cite{file:04_2_media_twophase}. This is the Van Genuchten parameter $m$, which is related to the pore-size distribution. Often $m = 1 - 1/n$ in VG models; an $m$ of 0.45 corresponds to $n \approx 1/(1-m) \approx 1.82$. Role: affects the shape of the retention curve $S(p_c)$ and the relative permeability function. A higher exponent (or higher n) makes the transition from residual to full saturation steeper. In the relative permeability model \texttt{RelativePermeabilityVanGenuchten}, this $m$ is used to compute $k_{\text{rel}}(S_e) = S_e^{1/2}[1 - (1 - S_e^{1/m})^m]^2$ (for Van Genuchten-Mualem) or similar \cite{file:04_2_media_twophase}. In the \texttt{SaturationVanGenuchten} model, it defines $p_c(S)$ relationship: $S_e = [1+ (p_c/p_b)^{n}]^{-m}$ with m=0.45. Thus, 0.45 is a moderate value, indicating a fairly gradual retention curve. These models directly couple into the Richards equation via the term $\phi \frac{\partial S}{\partial p_c}$ \cite{ogs:trm} and via $k_{\text{rel}}$ in Darcy's law \cite{ogs:trm}.
        \item \textbf{p\_b (Van Genuchten capillary pressure parameter)} – 1.0e+7 Pa (10 MPa) \cite{file:04_2_media_twophase}. This is the air-entry pressure (scale of capillary pressure at which saturation significantly drops). 10 MPa is very high, suggesting the material requires an extremely large suction to desaturate - effectively this implies the material will remain nearly fully saturated until an enormous tension is applied. In practical terms, this might simulate a nearly incompressible liquid in low-permeability rock (the domain likely stays close to saturation 1 for any realistic pressure changes). Role: in the retention curve, a higher $p_b$ means $S^l \approx 1$ for $p_c$ up to that magnitude. Given that pressure changes here are in the order of a few MPa at most, $p_c$ (which is -p in single-phase) will be far below 10 MPa in magnitude (since p is around 1.5 MPa, $p_c \approx -1.5$ MPa, effectively 0 suction), keeping $S^l \approx 1$ throughout. Thus, the model likely stays saturated or near-saturated. If unsaturation did occur, the formulas would use this $p_b$ to compute the actual saturation.
        \item \textbf{minimum\_relative\_permeability\_liquid} – 1e-25 (-) \cite{file:04_2_media_twophase}. This imposes a floor on the relative permeability for numerical stability (so that even at very low saturation, there is a tiny conductivity). Role: prevents complete shut-off of flow in dry conditions which could cause solver issues. In practice, because S stays high here, $k_{\text{rel}}$ likely stays near 1 and this minimum isn't active except at extreme dryness.
    \end{itemize}
\end{itemize}

These unsaturated parameters ensure the model has the capability for two-phase effects, though given the high $p_b$, the simulation might effectively treat the medium as fully saturated (liquid phase present everywhere) with negligible air influence. The Bishop's model exponent being 1 means effective stress uses $b(S)=S$ directly \cite{ogs:trm, doc:trm_gov_eq}, which at $S \approx 1$ means pore pressure fully enters the Terzaghi effective stress. If $S$ dropped (e.g. drying), then $b(S)$ would reduce, meaning the pore pressure has less effect on reducing effective stress (intuitively, if the pores empty of water, the remaining water pressure doesn't carry the load, the solid skeleton carries more).


\begin{thebibliography}{99}

\bibitem{ogs:trm}
OpenGeoSys official documentation – Thermo-Richards-Mechanics Process (TRM): governing equations and required properties.
\href{https://www.opengeosys.org/docs/processes/thermal-processes/trm/}{https://www.opengeosys.org/docs/processes/thermal-processes/trm/}

\bibitem{file:01_processes_TRM}
User's OGS project file: \texttt{01\_processes\_TRM.xml}. (processes setup)
\href{file://file_00000000c9007243b3805526bed5602f}{file://file\_00000000c9007243b3805526bed5602f}

\bibitem{file:05_time_loop_TRM}
User's OGS project file: \texttt{05\_time\_loop\_TRM.xml}.
\href{file://file_00000000025c724387276f9c32ff2436}{file://file\_00000000025c724387276f9c32ff2436}

\bibitem{doc:trm_gov_eq}
OpenGeoSys official documentation – Thermo-Richards-Mechanics Process (TRM): governing equations. (Cross-referenced from citation [2])
\href{https://www.opengeosys.org/docs/processes/thermal-processes/trm/}{https://www.opengeosys.org/docs/processes/thermal-processes/trm/}

\bibitem{file:04_media_TRM}
User's OGS project file: \texttt{04\_media\_TRM.xml}. (medium/phase property)
\href{file://file_00000000576472439d4d3944c05d5398}{file://file\_00000000576472439d4d3944c05d5398}

\bibitem{file:04_2_media_twophase}
User's OGS project file: \texttt{04\_2\_media\_twophase.xml}. (medium/phase property)
\href{file://file_00000000ad207243aaf25dfa733ce1be}{file://file\_00000000ad207243aaf25dfa733ce1be}

\bibitem{file:03_parameters_TRM}
User's OGS project file: \texttt{03\_parameters\_TRM.xml}. (defined parameters)
\href{file://file_000000000b947243a4e19a05569d7fb0}{file://file\_000000000b947243a4e19a05569d7fb0}

\bibitem{file:trm_medium_phase_prop}
User's OGS project file: (medium/phase property includes). (Cross-referenced from citation [22])

\bibitem{file:02_process_variables_TRM}
User's OGS project file: \texttt{02\_process\_variables\_TRM.xml}.
\href{file://file_00000000bf4871f4b7c1b2aad0def965}{file://file\_00000000bf4871f4b7c1b2aad0def965}

\bibitem{file:trm_proc_var_bc}
User's OGS project file: (process variables and boundary conditions). (Cross-referenced from citation [35, 47])

\bibitem{file:04_1_media_aqu_liq}
User's OGS project file: \texttt{04\_1\_media\_aqu\_liq.xml}.
\href{file://file_00000000a29c724385f8fd7a38b398f1}{file://file\_00000000a29c724385f8fd7a38b398f1}

\end{thebibliography}

\end{document}